\subsection{Repeated Failures and Concurrency}
\label{sec:sensitivity-results}
\begin{figure*}[tp]
\centering
\includegraphics[width=\textwidth]{figures/sensitivity.pdf}
\caption{Repeated-failure and concurrency sensitivity. Left: failure counts
1, 2, 4, and 8 at concurrency 16. Right: concurrency 8, 16, 24, and 32 with
two failures. Rows show absolute focus mean completion, focus mean and P95
gains, focus and victim workflow counts, Peer-HBM quota and pressure,
and selected Peer/Store prompt fractions. Quota is the time-weighted Peer
fraction; pressure is the fraction of control-window decisions observing
pressure. Points and whiskers summarize five
seeds; individual seed marks remain visible. Positive completion gains favor
AegisKV. Focus P95 is the within-seed maximum for these small focus sets.
The F2/C16 cell appears in both columns. Lines connect observed means;
they do not establish monotonicity. Natural targets and collateral victims
can differ between policies and across failure counts.}
\Description{Twelve plots showing the paired policies over failure count and
workflow concurrency, with resource and recovery diagnostics.}
\label{fig:sensitivity}
\end{figure*}


Figure~\ref{fig:sensitivity} summarizes all seven settings and five paired
seeds. At concurrency 16, mean focus WCT gains for one, two, four and eight
failures are $-8.82\%$, $-15.94\%$, $-0.15\%$ and $5.52\%$. With two failures,
concurrencies 8, 16, 24 and 32 give $-10.66\%$, $-15.94\%$, $6.32\%$ and
$-4.23\%$. Every focus-mean interval and every all-workflow-mean interval
crosses zero. The observed means support neither a consistently positive
gain nor a monotonic increase with failure count or concurrency.

Gains average seed-level percentages; they are not ratios of the absolute
means in row~(a), so their signs need not agree with that ratio. In F4/C16,
the focus P95 gain is 15.77\%, with interval [0.41\%, 31.14\%]. This isolated
positive comparison is exploratory and unadjusted; with four focus workflows,
P95 is the maximum rather than a resolved population-tail estimate.

Resource diagnostics also resist a simple pressure-to-benefit interpretation.
At two failures, increasing concurrency from 8 to 32 raises mean pressure
frequency from 3.42\% to 4.88\%, while mean Peer quota changes from 33.34\%
to 28.35\% and is not strictly decreasing at intermediate points. Across
the failure-count sweep, mean quota stays between 29.15\% and 30.12\%.
These averages describe controller exposure; the event-level audit separately
verifies its AIMD updates. Focus and collateral counts in row~(d) retain the
different realized victim populations rather than treating them as matched
samples. Selected Peer/Store fractions likewise describe recovery composition,
without asserting that tier overlap provides additional prefix coverage.
